%==============================================================
%  Cover letter for IEEE TIFS submission.
%  Compile with: pdflatex cover_letter
%  (Standalone; does not depend on the manuscript bundle.)
%==============================================================
\documentclass[11pt]{article}
\usepackage[margin=1in]{geometry}
\usepackage{parskip}
\usepackage[hidelinks]{hyperref}
\pagestyle{empty}

\begin{document}

\noindent Mihir Kavishwar\\
PixelRoot Project\\
\href{mailto:hello@mihirkavishwar.com}{hello@mihirkavishwar.com}\\[1em]

\noindent June 2026\\[1em]

\noindent To the Editor-in-Chief and Associate Editors,\\
\textit{IEEE Transactions on Information Forensics and Security}\\[1em]

\noindent Dear Editors,

I am pleased to submit the manuscript entitled \textbf{``PixelRoot:
Hardware-Anchored, Sensor-Level Media Provenance with Real-Time Blockchain
Notarization for Trustworthy Imaging''} for consideration as a regular paper in
the \textit{IEEE Transactions on Information Forensics and Security}.

The proliferation of photorealistic AI-generated imagery has outpaced
detector-based defenses, which remain locked in an adversarial arms race and
generalize poorly to unseen generators. Existing provenance schemes (e.g.,
C2PA) anchor trust in detachable metadata. This paper proposes an alternative:
moving the root of trust into the pixels themselves at the moment of capture,
combined with real-time blockchain notarization of a cryptographic commitment to
the frame. We believe this directly addresses core concerns of TIFS readers in
media forensics, watermarking, and content authenticity.

The principal contributions of the work are:
\begin{itemize}
  \item a hardware-rooted, dual-binding provenance architecture that combines an
        in-pixel keyed spread-spectrum signature (soft binding) with on-chain
        notarization of a per-frame hash (hard binding), organized into explicit
        assurance tiers with a software-ISP fallback for legacy devices;
  \item a concrete, tiered hardware ask, including a deep manufacturing
        feasibility analysis of per-pixel gain control (dual-conversion-gain
        pixels, coded-exposure pixels, and 3D-stacked digital-pixel sensors with
        Cu--Cu hybrid bonding), with engineering schematics grounded in the
        current fabrication literature;
  \item a formal security analysis (hard-binding soundness, carrier
        unpredictability, soft-binding false-accept bounds) and a
        privacy-preserving notarization design (salted commitments, selective
        disclosure, no raw PII on-chain); and
  \item an empirical evaluation conducted entirely on public images (the Kodak
        suite) with production codecs (\texttt{libjpeg}, \texttt{libwebp}) and a
        real Reed--Solomon error-correction layer, evaluated under a
        WAVES-aligned attack suite. We report imperceptibility (PSNR
        $\approx$~38~dB, SSIM~$\approx$~0.94), payload recovery (23/24 images
        decoded down to JPEG~Q10), a measured zero false-accept rate over 1224
        wrong-key/unmarked trials, and tamper localization at ROC-AUC~$0.97$, and
        we report the geometric-attack limitations honestly.
\end{itemize}

I confirm that this manuscript is original, has not been published previously,
and is not under consideration for publication elsewhere, in whole or in part.
The work is the sole author's. There are no conflicts of interest to disclose,
and the work received no external funding. The reference implementation and the
evaluation harness are released openly so that all reported numbers are
reproducible end-to-end.

I would be grateful for the opportunity to have this work reviewed and am happy
to provide any additional information the editors may require.

Thank you for your time and consideration.\\[1.5em]

\noindent Sincerely,\\[0.5em]
\noindent Mihir Kavishwar\\
PixelRoot Project

\end{document}
